% =====================================================================
%  Estimació de l'oxidació de carbohidrats i greixos a partir de la
%  freqüència cardíaca i de l'energia estimada per Garmin (article, en català)
%  Versió estable del 21/11/2026 (revisió: punts 1-77).
%  Compilar: pdflatex model-fc-energia.ca.tex (dues passades)
% =====================================================================
\documentclass[10pt,a4paper]{article}

\usepackage[utf8]{inputenc}
\usepackage[T1]{fontenc}
\usepackage[catalan]{babel}
\usepackage{lmodern}
\usepackage{microtype}
\usepackage[a4paper,margin=2.4cm]{geometry}
\usepackage{amsmath,amssymb}
\usepackage{icomma}
\usepackage{booktabs}
\usepackage{tabularx}
\usepackage{array}
\usepackage{caption}
\usepackage{float}
\usepackage{needspace}
\usepackage{enumitem}
\usepackage{tikz}
\usetikzlibrary{positioning,arrows.meta}
\usepackage{pgfplots}
\pgfplotsset{compat=1.17}
\usepackage[most]{tcolorbox}
\usepackage[numbers,sort&compress]{natbib}
\usepackage[hidelinks]{hyperref}

\captionsetup{font=small,labelfont=bf}
\renewcommand{\arraystretch}{1.25}
\newcolumntype{L}{>{\raggedright\arraybackslash}X}
\setlist{itemsep=2pt,topsep=4pt}

% Requadre verd per a les justificacions (el títol complet es passa com a paràmetre)
\newtcolorbox{requadre}[2][]{enhanced, colback=green!4, colframe=green!45!black,
  boxrule=0.6pt, arc=2pt, left=6pt, right=6pt, top=4pt, bottom=4pt,
  fonttitle=\bfseries\small, title={$\checkmark$\ #2}, #1}

% Requadre blau per a les conclusions d'una secció
\newtcolorbox{conclusio}{enhanced, colback=blue!4, colframe=blue!50!black,
  boxrule=0.6pt, arc=2pt, left=6pt, right=6pt, top=4pt, bottom=4pt,
  fonttitle=\bfseries\small, title={Conclusió}}

% Enllaç a la web del projecte (canviarà: modificar-lo només aquí)
\newcommand{\webprojecte}{\url{https://irazuzta.github.io/ChoFatMetabolic/}}

% --- Notació ---------------------------------------------------------
\newcommand{\sub}[1]{{\scriptscriptstyle\mathrm{#1}}}   % subíndexs petits
\newcommand{\CHO}{\sub{CHO}}
\newcommand{\FAT}{\sub{FAT}}
\newcommand{\tot}{\sub{total}}
\newcommand{\smax}{{\scriptscriptstyle\text{màx}}}
\newcommand{\cho}{\ensuremath{\mathit{CHO}}}               % ritme d'oxidació (g/min)
\newcommand{\fat}{\ensuremath{\mathit{FAT}}}
\newcommand{\VO}{\ensuremath{\mathit{VO}_2}}               % ritme (L/min)
\newcommand{\VCO}{\ensuremath{\mathit{VCO}_2}}
\newcommand{\VOmax}{\ensuremath{\mathit{VO}_{2\smax}}}
\newcommand{\RER}{\ensuremath{\mathit{RER}}}
\newcommand{\RERi}{\ensuremath{\overline{\mathit{RER}}}}
\newcommand{\FC}{\ensuremath{\mathit{FC}}}
\newcommand{\FCmax}{\ensuremath{\mathit{FC}_{\smax}}}
\newcommand{\Ion}{\ensuremath{I_{\mathrm{onset}}}}
\newcommand{\kcal}{\ensuremath{\mathrm{kcal}}}

\title{\textbf{Estimació de l'oxidació de carbohidrats i greixos\\
a partir de la freqüència cardíaca\\ i de l'energia estimada per Garmin}}
\author{Jordi Irazuzta\\
\small\href{mailto:irazuzta@gmail.com}{\texttt{irazuzta@gmail.com}}\\[4pt]
\small Projecte CHO/FAT Metabolic Pro --- \webprojecte}
\date{Esborrany --- novembre de 2026}

\begin{document}
\maketitle

\begin{abstract}
\noindent
Aquest treball estima el ritme d'oxidació de carbohidrats (CHO) i de greix (FAT) durant
l'exercici. Aquesta informació pot servir, per exemple, per decidir quants carbohidrats
ingerir. La calorimetria indirecta ho calcula amb les equacions de Jeukendrup i Wallis,
però necessita una anàlisi de gasos per mesurar el consum d'oxigen (\VO) i el quocient
d'intercanvi respiratori (\RER). Aquí ho estimem a partir d'un rellotge esportiu que mesura
la freqüència cardíaca (\FC) i dona una estimació de l'energia gastada; en concret, els
dispositius Garmin. El mètode parteix d'una condició: l'energia dels grams calculats ha de
ser exactament l'energia que dona Garmin. El \RER\ s'estima a partir de la \FC. El \VO\
s'obté dividint l'energia per l'energia que dona cada litre d'oxigen. Aquest valor es
calcula amb les mateixes equacions de Jeukendrup i Wallis i amb l'energia de cada gram de
carbohidrat (4,1~kcal) i de greix (9,3~kcal), uns valors que es justifiquen amb la
calorimetria clàssica. Així, la condició es compleix exactament. També s'analitzen
l'error en funció del \RER, com es passa del ritme al total, les fonts d'error i com aplicar-lo
en temps real.
\end{abstract}

\noindent\textbf{Paraules clau:} oxidació de substrats; calorimetria indirecta;
consum d'oxigen; quocient respiratori; despesa energètica; dispositius portables.

% =====================================================================
\section{Objectiu}\label{sec:objectiu}
% =====================================================================

Volem estimar, en cada instant $t$ de l'activitat, el \textbf{ritme d'oxidació} de
carbohidrats i de greix:
\[
  \cho(t),\ \fat(t)\qquad[\mathrm{g/min}],
\]
i, a partir d'aquests ritmes, el \textbf{total acumulat} de tota l'activitat, des de 0
fins a la durada total $T$:
\begin{equation}
  \cho_{\tot}(T)=\int_0^T\cho(t)\,dt,\qquad
  \fat_{\tot}(T)=\int_0^T\fat(t)\,dt .
  \label{eq:totals}
\end{equation}
La integral és el total acumulat: tot el que s'ha oxidat des de l'inici fins a $T$. La
secció~\ref{sec:total} explica com es calcula amb les dades de Garmin. Si es vol, aquesta informació pot servir per decidir
quants carbohidrats ingerir durant l'esforç \citep{jeukendrup2014}.

\paragraph{El plantejament.} La calorimetria indirecta calcula aquests ritmes a partir dels
gasos respiratoris (secció~\ref{sec:jw}), però necessita un analitzador de gasos. Aquí els
estimem amb els dispositius Garmin (rellotges i ciclocomputadors Edge), que mesuren la \FC\
i estimen l'energia gastada. Els passos són:
\begin{enumerate}[leftmargin=*,itemsep=0pt]
  \item De Garmin obtenim l'energia gastada, $E(t)$ (secció~\ref{sec:energia}).
  \item El \RER\ s'estima a partir de la \FC\ (secció~\ref{sec:rer}).
  \item A partir del \RER\ i de $E(t)$ podem estimar el \VO\ (secció~\ref{sec:vo2}).
  \item Amb el \VO\ i el \RER, les equacions de Jeukendrup i Wallis donen $\cho(t)$ i
        $\fat(t)$ (secció~\ref{sec:jw}).
\end{enumerate}

\paragraph{Per què aquest plantejament.} Una altra opció és estimar el consum d'oxigen a partir de la \FC, amb una relació fixa entre tots
dos \citep{swain1997,chofatfc2026}. Però llavors l'energia que surt dels grams calculats no
coincideix amb la que mostra Garmin: acostuma a sortir més alta. Garmin fa servir més
informació: a més de la \FC, el perfil de l'usuari i la relació de cada persona entre \FC\ i
consum d'oxigen \citep{firstbeat2007}.

Per això aquest treball parteix d'una \textbf{condició}: l'energia dels grams calculats ha
de ser exactament l'energia que dona Garmin. L'anomenem, amb una mica d'humor,
«principi de conservació de l'energia Garmin», per analogia amb el principi de la física:
aquí l'energia no la fixa la natura, sinó Garmin. El mètode no torna a calcular l'energia;
només la reparteix entre els dos combustibles. Aquesta condició és la que ho determina tot
(secció~\ref{sec:vo2}). La resta del plantejament i la notació segueixen el model que
estima el consum d'oxigen a partir de la \FC\ \citep{chofatfc2026}.

\paragraph{Organització.} La secció~\ref{sec:conceptes} introdueix els conceptes necessaris. Després es respon, per
ordre, a dues preguntes: \textbf{com es calculen} $\cho(t)$ i $\fat(t)$
(seccions~\ref{sec:jw}--\ref{sec:sensibilitat}) i \textbf{com es calculen els totals}
(secció~\ref{sec:total}). Al final hi ha un exemple, les fonts d'error i com
aplicar el mètode en temps real.

% =====================================================================
\section{Conceptes previs}\label{sec:conceptes}
% =====================================================================

\subsection{Consum d'oxigen (\texorpdfstring{$\mathit{VO}_2$}{VO2}) i producció de
CO\texorpdfstring{$_2$}{2} (\texorpdfstring{$\mathit{VCO}_2$}{VCO2})}\label{sec:gasos}
$\VO(t)$ és el volum d'oxigen que el cos consumeix per unitat de temps. Es mesura en
L/min (valor absolut) o en mL$\cdot$kg$^{-1}\cdot$min$^{-1}$ (relatiu al pes, per comparar
persones). $\VCO(t)$ és, de la mateixa manera, el volum de diòxid de carboni produït per
unitat de temps. Tots dos són la base de la \emph{calorimetria indirecta}: oxidar greix o
carbohidrats consumeix O$_2$ i produeix CO$_2$ en proporcions diferents, de manera que
mesurant $\VO$ i $\VCO$ es pot saber quin combustible s'està cremant sense mesurar-lo
directament.

\subsection{Consum màxim d'oxigen (\texorpdfstring{$\mathit{VO}_{2\max}$}{VO2max})}
$\VOmax$ és el valor més alt de $\VO$ que una persona pot assolir. Aquí només s'utilitza
per expressar intensitats (per exemple, ``al 63\,\% del $\VOmax$'') i per detectar valors
de $\VO(t)$ impossibles.

\subsection{Quocient d'intercanvi respiratori (\texorpdfstring{$\mathit{RER}$}{RER})}
\begin{equation}
  \RER(t)=\frac{\VCO(t)}{\VO(t)} .
  \label{eq:rer-def}
\end{equation}
És un sol nombre que indica la barreja de combustibles. Oxidar glucosa produeix tant
CO$_2$ com O$_2$ consumeix ($\RER=1$), mentre que oxidar greix en produeix menys
($\RER\approx0{,}70$). En exercici submàxim pren valors entre aquests dos extrems
\citep{peronnet1991}. És la variable central del repartiment entre carbohidrats i greix.

\subsection{Energia estimada pels dispositius Garmin}\label{sec:energia}
Els dispositius Garmin calculen l'energia de l'activitat amb un mètode de Firstbeat, una
empresa finlandesa especialitzada en l'anàlisi fisiològica a partir de la freqüència
cardíaca; la seva divisió de consum (Firstbeat Analytics) forma part de Garmin des del 2020. El mètode fa servir la \FC\ i el perfil de l'usuari, i té en compte la relació de cada persona entre
\FC\ i consum d'oxigen \citep{firstbeat2007}. Anomenem $E(t)$ el \emph{ritme} de despesa
d'energia, en kcal/min, i $E_{\tot}(t)$ l'energia \emph{acumulada} des de l'inici, en kcal:
\begin{equation}
  E_{\tot}(t)=\int_0^t E(s)\,ds .
  \label{eq:etotal}
\end{equation}
Garmin dona $E_{\tot}(t)$; el ritme $E(t)$ s'obté de les pujades d'aquest comptador
(secció~\ref{sec:total}). És l'energia \emph{total}: inclou també
la del metabolisme de repòs durant aquell temps (Garmin Connect la separa en calories en
repòs i calories actives). És la que cal, perquè les equacions de calorimetria compten tota
l'oxidació a partir de tot l'oxigen consumit. És una estimació, no una mesura: els estudis
sobre dispositius portables hi troben errors més grans que en la mesura de la \FC\
\citep{shcherbina2017,fuller2020}.

\needspace{10cm}
\subsection{Notació}
La taula~\ref{tab:notacio} recull els símbols del treball.

\begin{table}[H]
\centering
\caption{Notació.}
\label{tab:notacio}
\small
\begin{tabularx}{\textwidth}{@{}l L l@{}}
\toprule
Símbol & Significat & Unitats \\
\midrule
$\cho(t)$, $\fat(t)$ & Ritme d'oxidació de carbohidrats i de greix (el resultat) & g/min \\
$\cho_{\tot}(T)$, $\fat_{\tot}(T)$ & Total oxidat des de l'inici fins a $T$ & g \\
$\VO(t)$, $\VCO(t)$ & Consum d'O$_2$ i producció de CO$_2$ & L/min \\
$\RER(t)$ & Quocient d'intercanvi respiratori, $\VCO/\VO$ & --- \\
$E_{\tot}(t)$ & Energia acumulada des de l'inici (la dona Garmin) & kcal \\
$E(t)$ & Ritme de despesa d'energia (s'obté de $E_{\tot}$) & kcal/min \\
$g_{\CHO}(\RER)$, $g_{\FAT}(\RER)$ & Grams de substrat oxidats per cada litre d'O$_2$ & g/L \\
$e_{\CHO}$, $e_{\FAT}$ & Energia per gram de substrat oxidat; constants: 4,1 i 9,3 & kcal/g \\
$k(\RER)$ & Energia alliberada per cada litre d'O$_2$ & kcal/L \\
$\FC(t)$, $LT_1$, $LT_2$, $\FCmax$, $\Ion$ & Freqüència cardíaca i punts de referència del \RER & ppm \\
\bottomrule
\end{tabularx}
\end{table}

% =====================================================================
\section{El punt de partida: les equacions de Jeukendrup i Wallis}\label{sec:jw}
% =====================================================================

La relació entre l'oxigen consumit, el CO$_2$ produït i els grams de substrat oxidats es
coneix des de les equacions clàssiques de calorimetria indirecta de \citet{frayn1983}. Són
equacions «no proteiques»: suposen que tota l'energia surt dels carbohidrats i dels
greixos, i no tenen en compte l'oxidació de proteïnes, que durant l'exercici aporta poca
energia. Aquest treball fa servir la versió per a exercici de \citet{jeukendrup2005}:
\begin{align}
  \cho(t) &= 4{,}210\cdot\VCO(t)-2{,}962\cdot\VO(t), \label{eq:jw-cho}\\
  \fat(t) &= 1{,}695\cdot\VO(t)-1{,}701\cdot\VCO(t), \label{eq:jw-fat}
\end{align}
amb els gasos en L/min i el resultat en g/min.

\paragraph{Què cal per fer-les servir?} Les dues variables de la
secció~\ref{sec:gasos}: $\VO(t)$ i $\VCO(t)$. Cap de les dues es pot mesurar sense un
analitzador de gasos, i per tant cal estimar-les.

\subsection{Reescriure les equacions amb el RER}\label{sec:jw-rer}
En lloc d'estimar $\VO$ i $\VCO$ per separat, és més pràctic estimar $\VO$ i el
\RER~\eqref{eq:rer-def}. Substituint $\VCO(t)=\RER(t)\cdot\VO(t)$:
\begin{align}
  \cho(t) &= \VO(t)\cdot\bigl(4{,}210\cdot\RER(t)-2{,}962\bigr), \label{eq:jw-cho-rer}\\
  \fat(t) &= \VO(t)\cdot\bigl(1{,}695-1{,}701\cdot\RER(t)\bigr). \label{eq:jw-fat-rer}
\end{align}

\paragraph{Els extrems.} Fora del rang 0,704--0,9965 un dels dos parèntesis surt negatiu;
per això es posa a zero ($\cho=0$ per sota, $\fat=0$ per sobre).

\needspace{12cm}
\paragraph{Ara calen dues estimacions:} $\RER(t)$ (secció~\ref{sec:rer}) i $\VO(t)$
(secció~\ref{sec:vo2}). El \RER\ va primer, perquè per estimar el $\VO$ cal conèixer-lo.
La figura~\ref{fig:esquema} resumeix el camí.

\begin{figure}[H]
\centering
\begin{tikzpicture}[
  node distance=6mm and 7mm,
  caixa/.style={draw, rounded corners=2pt, align=center, font=\small, minimum height=9mm, inner sep=4pt},
  dada/.style={caixa, fill=gray!12},
  pas/.style={caixa, fill=blue!7},
  clau/.style={caixa, fill=orange!20, very thick},
  res/.style={caixa, fill=green!12},
  fl/.style={-{Stealth[length=2.2mm]}, thick}]
\node[dada] (fc) {$\FC(t)$\\(Garmin)};
\node[pas, right=of fc] (rer) {$\RER(t)$\\ \S\ref{sec:rer}};
\node[pas, right=of rer] (g) {$g_{\CHO},\ g_{\FAT}$\\ \S\ref{sec:g}};
\node[clau, right=of g] (k) {$k(\RER)$\\ \S\ref{sec:k}};
\node[pas, above=of k] (e) {$e_{\CHO},\ e_{\FAT}$ \ \S\ref{sec:e}};
\node[dada, below=of fc, yshift=-4mm] (en) {$E_{\tot}(t)$ (Garmin)\\ $\to$ ritme $E(t)$};
\node[pas, below=of k, yshift=-4mm] (v) {$\VO(t)=E(t)/k$\\ \S\ref{sec:vo2-calcul}};
\node[res, below=of v] (m) {$\cho(t),\ \fat(t)$\\ \S\ref{sec:grams}};
\node[res, left=12mm of m] (t) {$\cho_{\tot},\ \fat_{\tot}$\\ \S\ref{sec:total}};
\draw[fl] (fc) -- (rer);
\draw[fl] (rer) -- (g);
\draw[fl] (g) -- (k);
\draw[fl] (e) -- (k);
\draw[fl] (k) -- (v);
\draw[fl] (en) -- (v);
\draw[fl] (v) -- (m);
\draw[fl] (m) -- (t);
\end{tikzpicture}
\caption{Esquema del mètode. La \FC\ dona el $\RER(t)$, i el \RER\ els grams oxidats per
litre d'oxigen, $g_{\CHO}$ i $g_{\FAT}$. Amb l'energia per gram de cada substrat, $e_{\CHO}$ i
$e_{\FAT}$, s'obté l'energia per litre d'oxigen, $k$, que converteix el ritme d'energia que
s'obté de Garmin, $E(t)$, en consum d'oxigen, $\VO(t)$. Amb el consum d'oxigen i el \RER\
s'obté el ritme d'oxidació, $\cho(t)$ i $\fat(t)$, i a partir d'aquí els totals,
$\cho_{\tot}$ i $\fat_{\tot}$.}
\label{fig:esquema}
\end{figure}

% =====================================================================
\section{Estimar \texorpdfstring{$\mathit{RER}(t)$}{RER(t)}: llindars metabòlics}\label{sec:rer}
% =====================================================================

El \RER\ puja amb la intensitat, sobretot al voltant dels llindars metabòlics: el primer
($LT_1$, llindar aeròbic) i el segon ($LT_2$, llindar de lactat o anaeròbic). S'estima a
partir de la \FC\ tal com es descriu a \citep{chofatfc2026}, on hi ha la justificació
completa. Aquí només se'n dona el resultat.

\subsection{Model d'interpolació}
Es fixen quatre punts de \FC\ ($\Ion$, $LT_1$, $LT_2$ i $\FCmax$) i un valor de \RER\ per a
cadascun (taula~\ref{tab:ancoratges}). Entre dos punts consecutius interpolem linealment.
Com que tenim 4 punts de referència, obtenim una funció definida a trossos amb 5 parts:
constant fins a $\Ion$, tres trams rectes entre punts, i constant a partir de $\FCmax$:
\begin{equation}
\RER(t)=
\begin{cases}
R_{\mathrm{basal}} & \FC(t)\le\Ion,\\[4pt]
R_{\mathrm{basal}}+\dfrac{\FC(t)-\Ion}{LT_1-\Ion}\cdot\bigl(R_{LT_1}-R_{\mathrm{basal}}\bigr) & \Ion<\FC(t)\le LT_1,\\[10pt]
R_{LT_1}+\dfrac{\FC(t)-LT_1}{LT_2-LT_1}\cdot\bigl(R_{LT_2}-R_{LT_1}\bigr) & LT_1<\FC(t)\le LT_2,\\[10pt]
R_{LT_2}+\dfrac{\FC(t)-LT_2}{\FCmax-LT_2}\cdot\bigl(R_{\smax}-R_{LT_2}\bigr) & LT_2<\FC(t)\le\FCmax,\\[6pt]
R_{\smax} & \FC(t)>\FCmax.
\end{cases}
\label{eq:rer}
\end{equation}

\begin{table}[H]
\centering
\caption{Valors del \RER\ als quatre punts de referència \citep{chofatfc2026}.}
\label{tab:ancoratges}
\small
\begin{tabular}{@{}llc@{}}
\toprule
Punt de \FC & Nom & \RER \\
\midrule
Fins a $\Ion$ (esforç suau) & $R_{\mathrm{basal}}$ & 0,86 \\
$LT_1$ & $R_{LT_1}$ & 0,91 \\
$LT_2$ & $R_{LT_2}$ & 0,98 \\
$\FCmax$ i per sobre & $R_{\smax}$ & 1,00 \\
\bottomrule
\end{tabular}
\end{table}

Per sobre de $LT_2$ el \RER\ mesurat pot passar d'1, però aquest excés no és oxidació de
carbohidrats \citep{chofatfc2026}. Per això el sostre és 1,00, on el greix val zero
(secció~\ref{sec:jw-rer}).

\subsection{Paràmetres}\label{sec:parametres}
Si l'esportista s'ha fet una prova de lactat o d'esforç amb gasos, es fan servir els seus
valors reals. Si no, s'estimen com a \citep{chofatfc2026}: $LT_1\approx0{,}89\cdot LT_2$ i
$\Ion\approx0{,}70\cdot LT_2$.

Per calcular els grams no calen el pes ni el $\VOmax$: ja hi són dins de l'energia que
calcula Garmin. Però aquesta energia depèn del perfil de l'usuari al dispositiu (pes,
edat, \FC\ màxima, $\VOmax$), i per tant cal tenir-lo al dia.

El \RER\ pesa poc en el consum d'oxigen i molt en el repartiment entre carbohidrats i greix
(secció~\ref{sec:sensibilitat}). Per això el que més importa és que els llindars, sobretot
$LT_2$, estiguin ben posats.

% =====================================================================
\section{Estimar \texorpdfstring{$\mathit{VO}_2(t)$}{VO2(t)} a partir de l'energia de Garmin}\label{sec:vo2}
% =====================================================================

\paragraph{Què tenim.} El ritme de despesa d'energia $E(t)$, que s'obté a partir de
l'energia acumulada que dona Garmin, $E_{\tot}(t)$ (seccions~\ref{sec:energia}
i~\ref{sec:total}), i el $\RER(t)$, que s'estima amb la \FC\ (secció~\ref{sec:rer}).

\paragraph{Què volem.} El consum d'oxigen $\VO(t)$, amb la condició de la
secció~\ref{sec:objectiu}: l'energia dels grams que en surtin ha de ser exactament $E(t)$.

\paragraph{Com ho fem.} Si multipliquem les kcal que allibera cada litre d'oxigen en
l'instant $t$, $k\bigl(\RER(t)\bigr)$ (depèn de la barreja de combustibles d'aquell moment),
pels litres d'oxigen que es consumeixen per minut en aquell instant, $\VO(t)$, obtenim el
ritme de despesa d'energia en l'instant $t$, $E(t)$, en kcal/min:
\begin{equation}
  E(t)=k\bigl(\RER(t)\bigr)\cdot\VO(t) .
  \label{eq:energia-k}
\end{equation}
L'energia $E(t)$ la podem obtenir del que dona Garmin. Per tant, si estimem $k$ (l'energia per litre), només
cal aïllar el consum d'oxigen:
\begin{equation}
  \VO(t)=\frac{E(t)}{k\bigl(\RER(t)\bigr)} .
  \label{eq:vo2}
\end{equation}
$k$ és la peça clau. Es construeix en tres passos: els grams que crema
cada litre d'oxigen (secció~\ref{sec:g}), l'energia de cada gram (secció~\ref{sec:e}) i,
multiplicant, l'energia de cada litre (secció~\ref{sec:k}).

\subsection{Grams de cada substrat per litre d'oxigen}\label{sec:g}
\paragraph{Per què.} Ens interessa saber quants grams de cada substrat es cremen per
\emph{cada litre} d'oxigen. Així, quan sapiguem quants litres per minut es consumeixen,
només caldrà multiplicar.

\paragraph{Com.} Partim de les equacions de Jeukendrup i Wallis amb el
\RER~\eqref{eq:jw-cho-rer}--\eqref{eq:jw-fat-rer}:
\[
  \cho(t)=\VO(t)\cdot\bigl(4{,}210\cdot\RER(t)-2{,}962\bigr),\qquad
  \fat(t)=\VO(t)\cdot\bigl(1{,}695-1{,}701\cdot\RER(t)\bigr).
\]
Dividim els dos costats per $\VO(t)$:
\[
  \frac{\cho(t)}{\VO(t)}=4{,}210\cdot\RER(t)-2{,}962,\qquad
  \frac{\fat(t)}{\VO(t)}=1{,}695-1{,}701\cdot\RER(t) .
\]
A l'esquerra hi ha grams per minut dividits per litres per minut, és a dir, grams per
litre: (g/min)\,/\,(L/min) = g/L. Són els grams que crema cada
litre d'oxigen, i els anomenem:
\begin{align}
  g_{\CHO}(\RER) &= 4{,}210\cdot\RER-2{,}962 \qquad[\mathrm{g/L}], \label{eq:gc}\\
  g_{\FAT}(\RER) &= 1{,}695-1{,}701\cdot\RER \qquad[\mathrm{g/L}]. \label{eq:gf}
\end{align}
Només depenen del \RER. Per exemple, amb $\RER=0{,}86$ cada litre d'oxigen crema 0,66~g
de carbohidrats i 0,23~g de greix. Als extrems es posen a zero, com a la
secció~\ref{sec:jw-rer}.

\paragraph{\RER\ o $\RER(t)$.} Quan definim una funció del \RER, com $g_{\CHO}(\RER)$ o
$k(\RER)$, escrivim només \RER: és la variable de la funció i pot prendre qualsevol valor.
Quan l'apliquem en un instant, escrivim $\RER(t)$: per exemple, $g_{\CHO}\bigl(\RER(t)\bigr)$
són els grams per litre en l'instant $t$.

\subsection{Energia per gram de cada substrat}\label{sec:e}
\paragraph{Què són.} Oxidar completament un gram de substrat allibera una quantitat fixa
d'energia, que depèn només de la seva composició química: la seva \emph{densitat
energètica}. El mètode fa servir dues constants:
\begin{equation}
  e_{\CHO}=4{,}1~\kcal/\mathrm{g},\qquad e_{\FAT}=9{,}3~\kcal/\mathrm{g}.
  \label{eq:e}
\end{equation}

\paragraph{D'on surten.} La taula~\ref{tab:densitats} recull els valors de referència.
Els factors d'Atwater (4 i 9~kcal/g) són els que s'usen en nutrició per als aliments
\citep{atwater1900}. En calorimetria es fan servir els valors del substrat concret que
s'oxida \citep{livesey1988}: durant l'exercici, el carbohidrat oxidat és sobretot
glicogen i el greix, triacilglicerols. Els valors triats es troben en aquest rang, prop
del glicogen i dels triacilglicerols, que són precisament els substrats per als quals
estan pensades les equacions de Jeukendrup i Wallis.

\begin{table}[H]
\centering
\caption{Energia alliberada per gram de substrat oxidat.}
\label{tab:densitats}
\small
\begin{tabular}{@{}lcc@{}}
\toprule
Font & Carbohidrats (kcal/g) & Greixos (kcal/g) \\
\midrule
Factors d'Atwater (aliments) \citep{atwater1900} & 4 & 9 \\
Glucosa / glicogen oxidats \citep{livesey1988} & $\approx3{,}7$ / $\approx4{,}2$ & --- \\
Triacilglicerols oxidats \citep{livesey1988} & --- & $\approx9{,}4$ \\
\textbf{Aquest treball} ($e_{\CHO}$, $e_{\FAT}$) & \textbf{4,1} & \textbf{9,3} \\
\bottomrule
\end{tabular}
\end{table}

\paragraph{Per què són vàlids.} A més d'encaixar amb la literatura, els dos valors es
justifiquen amb dues comprovacions, que es fan a la secció~\ref{sec:k} i es marquen amb un
requadre verd:
\begin{enumerate}[label=(\alph*)]
  \item reprodueixen la taula de calorimetria de \citet{peronnet1991} als dos extrems;
  \item l'energia per litre d'oxigen que en resulta coincideix amb la fórmula clàssica de
        \citet{weir1949} amb menys d'un 1,5\,\% de diferència.
\end{enumerate}
A més, a la secció~\ref{sec:grams} es comprova que es compleix la condició de partida:
l'energia dels grams calculats és exactament l'energia que dona Garmin.

\subsection{Energia per litre d'oxigen, \texorpdfstring{$k(\mathit{RER})$}{k(RER)}}\label{sec:k}
És la peça clau del mètode. El raonament és directe:
\begin{itemize}
  \item cada litre d'oxigen oxida $g_{\CHO}$ grams de carbohidrat i $g_{\FAT}$ grams de
        greix (secció~\ref{sec:g});
  \item cada gram allibera $e_{\CHO}$ o $e_{\FAT}$ kilocalories (secció~\ref{sec:e});
  \item per tant, cada litre d'oxigen allibera
\end{itemize}
\begin{equation}
  k(\RER)=e_{\CHO}\cdot g_{\CHO}(\RER)+e_{\FAT}\cdot g_{\FAT}(\RER)\qquad[\kcal/\mathrm{L}].
  \label{eq:k}
\end{equation}

\paragraph{Una fórmula més simple.} Entre $\RER\approx0{,}704$ i $0{,}9965$, on es cremen
tots dos substrats, es poden substituir~\eqref{eq:gc}, \eqref{eq:gf} i~\eqref{eq:e}
a~\eqref{eq:k} i agrupar termes:
\begin{align}
  k(\RER) &= 4{,}1\cdot(4{,}210\cdot\RER-2{,}962)+9{,}3\cdot(1{,}695-1{,}701\cdot\RER) \notag\\
          &= 17{,}261\cdot\RER-12{,}144+15{,}764-15{,}819\cdot\RER \notag\\
          &= (17{,}261-15{,}819)\cdot\RER+(15{,}764-12{,}144) \notag\\
          &= 3{,}620+1{,}442\cdot\RER . \label{eq:klin}
\end{align}
Amb aquesta forma es veu de seguida que $k$ varia poc: va de 4,63~kcal/L (amb
$\RER=0{,}704$) a 5,06~kcal/L (amb $\RER=0{,}9965$). Per tant, energia i oxigen són gairebé
proporcionals. Que variï poc no vol dir que sigui constant: es calcula amb el \RER\ de
cada instant.

\begin{requadre}{Justificació de $e_{\CHO}$ i $e_{\FAT}$ --- (a) Calorimetria clàssica}
\textbf{Què es comprova:} que $e_{\CHO}$ i $e_{\FAT}$ donen la mateixa energia per litre
d'oxigen que la calorimetria clàssica quan només es crema un substrat.

\smallskip
Segons la taula de \citet{peronnet1991}, un litre d'O$_2$ allibera 5,047~kcal quan només
es crema carbohidrat i 4,686~kcal quan només es crema greix. En aquests extrems, cada litre
crema $g_{\CHO}(1{,}00)=1{,}248$~g de carbohidrat o $g_{\FAT}(0{,}70)=0{,}504$~g de greix.
Dividint l'energia pels grams:
\[
  \frac{5{,}047}{1{,}248}\approx4{,}05~\kcal/\mathrm{g},\qquad
  \frac{4{,}686}{0{,}504}\approx9{,}30~\kcal/\mathrm{g}.
\]
\textbf{Resultat:} són els valors de~\eqref{eq:e}.
\end{requadre}

\begin{requadre}{Justificació de $e_{\CHO}$ i $e_{\FAT}$ --- (b) Fórmula de Weir}
\textbf{Què es comprova:} que l'energia per litre d'oxigen $k(\RER)$ coincideix amb la
fórmula clàssica de \citet{weir1949}. Aquesta fórmula no s'utilitza en el mètode: només
serveix per comprovar $k$.

\smallskip
\textbf{D'on surt la fórmula de Weir en funció del \RER.} Weir relaciona gasos i energia
(sense el terme de nitrogen de l'orina):
\[
  E(t)=3{,}941\cdot\VO(t)+1{,}106\cdot\VCO(t).
\]
Com a la secció~\ref{sec:jw-rer}, substituïm $\VCO(t)=\RER(t)\cdot\VO(t)$:
\[
  E(t)=3{,}941\cdot\VO(t)+1{,}106\cdot\RER(t)\cdot\VO(t)=\VO(t)\cdot\bigl(3{,}941+1{,}106\cdot\RER(t)\bigr).
\]
Dividim per $\VO(t)$ per tenir l'energia de cada litre, que és el mateix que és $k$:
\[
  \frac{E(t)}{\VO(t)}=3{,}941+1{,}106\cdot\RER(t)\qquad[\kcal/\mathrm{L}].
\]
La comparem amb la nostra, $k(\RER)=3{,}620+1{,}442\cdot\RER$~\eqref{eq:klin}: són dues
rectes molt properes.

\medskip
\noindent
\begin{minipage}[c]{0.44\linewidth}
\centering
\footnotesize
\begin{tabular}{@{}cccc@{}}
\toprule
\RER & $k$ & Weir & diferència \\
\midrule
0,70 & 4,690 & 4,715 & $-0{,}5\,\%$ \\
0,80 & 4,773 & 4,826 & $-1{,}1\,\%$ \\
0,86 & 4,859 & 4,892 & $-0{,}7\,\%$ \\
0,91 & 4,931 & 4,947 & $-0{,}3\,\%$ \\
0,95 & 4,989 & 4,992 & $-0{,}1\,\%$ \\
0,98 & 5,032 & 5,025 & $+0{,}1\,\%$ \\
1,00 & 5,117 & 5,047 & $+1{,}4\,\%$ \\
\bottomrule
\end{tabular}

\smallskip
$k$ i Weir en kcal per litre d'O$_2$.
\end{minipage}\hfill
\begin{minipage}[c]{0.53\linewidth}
\centering
\begin{tikzpicture}
\begin{axis}[
  width=\linewidth, height=4.6cm,
  xmin=0.70, xmax=1.00, xlabel={\RER},
  ylabel={kcal/L d'O$_2$}, ymin=4.6, ymax=5.15, grid=major, grid style={gray!25},
  xticklabel style={/pgf/number format/fixed,/pgf/number format/precision=2},
  yticklabel style={/pgf/number format/fixed,/pgf/number format/precision=1},
  tick label style={font=\scriptsize}, label style={font=\scriptsize},
  /pgf/number format/use comma,
  legend style={font=\scriptsize, at={(0.03,0.97)}, anchor=north west, draw=none, fill=white, fill opacity=0.8}]
\addplot[thick, black] coordinates {
 (0.7036,4.6337) (0.72,4.6573) (0.74,4.6862) (0.76,4.7150) (0.78,4.7438) (0.80,4.7727)
 (0.82,4.8015) (0.84,4.8303) (0.86,4.8592) (0.88,4.8880) (0.90,4.9168) (0.92,4.9457)
 (0.94,4.9745) (0.96,5.0033) (0.98,5.0322) (0.9965,5.0564) (1.00,5.1168)};
\addlegendentry{$k(\RER)$}
\addplot[thick, dashed, black] coordinates {(0.70,4.7152) (1.00,5.0470)};
\addlegendentry{Weir}
\end{axis}
\end{tikzpicture}
\end{minipage}

\medskip
\textbf{Resultat:} la diferència és inferior a l'1,5\,\% en tot el rang. El petit salt del
final de la corba, entre $\RER=0{,}9965$ i $1{,}00$, és perquè en aquest tram el greix ja
val zero (secció~\ref{sec:jw-rer}) i $k$ deixa de ser la recta: només queda
$e_{\CHO}\cdot g_{\CHO}$. Per això la fila d'1,00 dona $+1{,}4\,\%$.
\end{requadre}

\subsection{Consum d'oxigen a partir de l'energia}\label{sec:vo2-calcul}
Amb $k$ ja calculada, s'aplica~\eqref{eq:vo2}. Les unitats quadren:
\[
  \VO(t)=\frac{E(t)}{k\bigl(\RER(t)\bigr)}\qquad
  \Bigl[\frac{\kcal/\mathrm{min}}{\kcal/\mathrm{L}}=\frac{\mathrm{L}}{\mathrm{min}}\Bigr].
\]
Per exemple, si es gasten 12~kcal/min i cada litre n'allibera 4,94, es consumeixen
$12/4{,}94=2{,}43$~L/min d'oxigen. Aquest valor es pot comparar amb el $\VOmax$ per
detectar valors impossibles.

\subsection{Ritme d'oxidació de cada substrat}\label{sec:grams}
Tornem a les equacions de Jeukendrup i Wallis amb el
\RER~\eqref{eq:jw-cho-rer}--\eqref{eq:jw-fat-rer}:
\[
  \cho(t)=\VO(t)\cdot\bigl(4{,}210\cdot\RER(t)-2{,}962\bigr).
\]
Hi posem el consum d'oxigen que hem obtingut, $\VO(t)=E(t)/k\bigl(\RER(t)\bigr)$~\eqref{eq:vo2}:
\[
  \cho(t)=\frac{E(t)}{k\bigl(\RER(t)\bigr)}\cdot\bigl(4{,}210\cdot\RER(t)-2{,}962\bigr).
\]
I amb la forma simple de $k$~\eqref{eq:klin} queda una fórmula que només depèn de l'energia
i del \RER:
\begin{align}
  \cho(t) &= E(t)\cdot\frac{4{,}210\cdot\RER(t)-2{,}962}{3{,}620+1{,}442\cdot\RER(t)}, \label{eq:cho-final}\\[4pt]
  \fat(t) &= E(t)\cdot\frac{1{,}695-1{,}701\cdot\RER(t)}{3{,}620+1{,}442\cdot\RER(t)}. \label{eq:fat-final}
\end{align}

Aquestes fórmules valen entre $\RER\approx0{,}704$ i $0{,}9965$. Fora d'aquest rang, els
extrems es posen a zero (secció~\ref{sec:jw-rer}) i es fa servir~\eqref{eq:k}. La quantitat
d'energia ve de Garmin, i el \RER\ determina com es reparteix.

\begin{requadre}{«Principi de conservació de l'energia Garmin»}
\textbf{Què es comprova:} que es compleix la condició amb què s'ha construït el mètode
(secció~\ref{sec:objectiu}): l'energia dels grams calculats és exactament l'energia que ha
donat Garmin.

\smallskip
Multiplicant cada ritme per la seva energia per gram i sumant:
\begin{equation}
\begin{aligned}
  e_{\CHO}\cdot\cho(t)+e_{\FAT}\cdot\fat(t)
  &=\VO(t)\cdot\Bigl(e_{\CHO}\cdot g_{\CHO}\bigl(\RER(t)\bigr)+e_{\FAT}\cdot g_{\FAT}\bigl(\RER(t)\bigr)\Bigr)\\
  &=\frac{E(t)}{k\bigl(\RER(t)\bigr)}\cdot k\bigl(\RER(t)\bigr)=E(t) .
\end{aligned}
  \label{eq:conservacio}
\end{equation}
\textbf{Resultat:} es compleix exactament, sigui quin sigui el \RER, també als extrems on
un dels dos valors $g$ és zero.

\smallskip
Passa perquè $k$ s'ha calculat amb els mateixos nombres que les equacions de Jeukendrup i
Wallis. No justifica els valors 4,1 i 9,3 (es compliria amb qualsevol parell de valors, si
$k$ es calcula amb els mateixos), però garanteix que la quantitat d'energia sempre és la de
Garmin, i que el mètode només la reparteix entre els dos substrats. Amb una $k$ fixa o
amb la de Weir, el consum d'oxigen seria pràcticament el mateix, però l'energia dels grams
no quadraria amb la de Garmin (diferències de fins a un 4\,\% i un 1,5\,\%).
\end{requadre}

% =====================================================================
\section{Estudi de l'error en funció del RER}\label{sec:sensibilitat}
% =====================================================================

El \RER\ s'estima a partir de la \FC, i el valor real de cada persona pot ser una mica
diferent (secció~\ref{sec:limitacions}). Aquesta secció estudia com canvien els grams
calculats quan el \RER\ varia una quantitat $\varepsilon$, amb l'energia $E$ fixada (la de
Garmin). Per simplificar, no escrivim la dependència de $t$.

Les fórmules~\eqref{eq:cho-final}--\eqref{eq:fat-final} tenen la forma
\[
  \cho(\RER)=E\cdot\frac{a\cdot\RER-b}{k(\RER)},\qquad
  \fat(\RER)=E\cdot\frac{a'-b'\cdot\RER}{k(\RER)},\qquad
  k(\RER)=c+d\cdot\RER,
\]
amb $a=4{,}210$, $b=2{,}962$, $a'=1{,}695$, $b'=1{,}701$, $c=3{,}620$ i $d=1{,}442$.

\paragraph{Variació exacta.} Si el \RER\ passa de $\RER$ a $\RER+\varepsilon$:
\begin{equation}
\begin{aligned}
  \Delta\cho&=\cho(\RER+\varepsilon)-\cho(\RER)
  =E\cdot\left(\frac{a\,(\RER+\varepsilon)-b}{k(\RER+\varepsilon)}-\frac{a\cdot\RER-b}{k(\RER)}\right)\\[4pt]
  &=\frac{(a\,c+b\,d)\cdot E\cdot\varepsilon}{k(\RER)\cdot k(\RER+\varepsilon)}
  =\frac{19{,}51\cdot E\cdot\varepsilon}{k(\RER)\cdot k(\RER+\varepsilon)} ,
\end{aligned}
\label{eq:dcho}
\end{equation}
i, de la mateixa manera,
\begin{equation}
  \Delta\fat=-\frac{(b'c+a'd)\cdot E\cdot\varepsilon}{k(\RER)\cdot k(\RER+\varepsilon)}
  =-\frac{8{,}60\cdot E\cdot\varepsilon}{k(\RER)\cdot k(\RER+\varepsilon)} .
\label{eq:dfat}
\end{equation}
Al numerador, els termes amb \RER\ se cancel·len. La variació és proporcional a
$\varepsilon$ i a l'energia $E$, i el \RER\ només hi apareix a través de $k$, que varia poc.

L'energia total no canvia: $4{,}1\cdot19{,}51\approx9{,}3\cdot8{,}60\approx80$, de manera
que l'energia que guanyen els carbohidrats és la que perd el greix. És el «principi de
conservació de l'energia Garmin»~\eqref{eq:conservacio}.

\paragraph{Aproximació.} Per a variacions petites, $k(\RER+\varepsilon)\approx k(\RER)$:
\[
  \Delta\cho\approx\frac{19{,}51}{k(\RER)^2}\cdot E\cdot\varepsilon,\qquad
  \Delta\fat\approx-\frac{8{,}60}{k(\RER)^2}\cdot E\cdot\varepsilon .
\]
L'error d'aquest pas és molt petit: menys de l'1,5\,\% fins i tot amb $\varepsilon=0{,}05$.

\paragraph{Fita.} Dins del rang vàlid del \RER\ (0,704--0,9965), $k$ va de 4,63 a 5,06
(taula~\ref{tab:sens}).

\begin{table}[H]
\centering
\caption{Factors de~\eqref{eq:dcho}--\eqref{eq:dfat} aproximats, als extrems i al mig del
rang vàlid del \RER.}
\label{tab:sens}
\small
\begin{tabular}{@{}ccccc@{}}
\toprule
\RER & $k$ (kcal/L) & $k^2$ & $\dfrac{19{,}51}{k^2}$ & $\dfrac{8{,}60}{k^2}$ \\[6pt]
\midrule
0,704  & 4,63 & 21,5 & 0,91 & 0,40 \\
0,85   & 4,85 & 23,5 & 0,83 & 0,37 \\
0,9965 & 5,06 & 25,6 & 0,76 & 0,34 \\
\bottomrule
\end{tabular}
\end{table}

\noindent El factor dels carbohidrats queda sempre entre 0,76 i 0,91. Prenent el valor del
mig, 0,83, l'error és com a màxim d'un $\pm9$\,\%. Amb $E$ en kcal/min, i multiplicant per
60 per passar a g/h:
\begin{equation}
  \Delta\cho\approx0{,}83\cdot E\cdot\varepsilon~[\mathrm{g/min}]
  \approx50\cdot E\cdot\varepsilon~[\mathrm{g/h}],\qquad
  \Delta\fat\approx-0{,}37\cdot E\cdot\varepsilon~[\mathrm{g/min}]
  \approx-22\cdot E\cdot\varepsilon~[\mathrm{g/h}].
\label{eq:dsimple}
\end{equation}

\paragraph{Exemple.} Amb $E=12$~kcal/min i $\varepsilon=0{,}01$:
$\Delta\cho\approx50\cdot12\cdot0{,}01=6$~g/h i $\Delta\fat\approx-22\cdot12\cdot0{,}01
\approx-2{,}6$~g/h. Amb la fórmula exacta~\eqref{eq:dcho} surt 5,9~g/h amb $\RER=0{,}86$ i
5,6~g/h amb $\RER=0{,}95$. Si el \RER\ real d'una persona és 0,04 diferent de l'estimat, la
diferència és d'uns 24~g/h de carbohidrats.

\begin{conclusio}
\begin{itemize}[leftmargin=*]
  \item La variació dels grams és proporcional a la variació del \RER\ i a l'energia, i
        gairebé no depèn del valor del \RER.
  \item \textbf{Regla pràctica: $\Delta\cho\approx50\cdot E\cdot\varepsilon$~g/h i
        $\Delta\fat\approx-22\cdot E\cdot\varepsilon$~g/h} ($E$ en kcal/min). A 12~kcal/min,
        cada 0,01 de \RER\ són uns 6~g/h de carbohidrats, i uns 2,6~g/h de greix en sentit
        contrari.
  \item L'energia total no canvia: el \RER\ només determina com es reparteix.
  \item Per això el que més importa és tenir els llindars, sobretot $LT_2$, ben posats.
\end{itemize}
\end{conclusio}

% =====================================================================
\section{Càlcul de \texorpdfstring{$\mathit{CHO}_{\mathrm{total}}$ i $\mathit{FAT}_{\mathrm{total}}$}{CHO total i FAT total}}\label{sec:total}
% =====================================================================

L'objectiu~\eqref{eq:totals} demana el total acumulat, és a dir, la integral dels ritmes al
llarg de l'activitat. Aquesta secció explica com es calcula amb les dades que dona Garmin.

\subsection{Garmin dona l'energia acumulada}
Garmin no dona el ritme $E(t)$, sinó l'energia acumulada $E_{\tot}(t)$~\eqref{eq:etotal},
com un nombre enter de quilocalories que puja d'una en una \citep{garminciq}. En activitats
reals amb un rellotge Garmin s'ha comprovat que:
\begin{itemize}
  \item el comptador és l'energia \emph{total} (repòs inclòs): el seu valor final coincideix
        amb el total de l'activitat;
  \item puja d'1~kcal cada vegada, a intervals que depenen de la intensitat (cada
        5--7~s corrent);
  \item no avança amb el cronòmetre aturat.
\end{itemize}
Es pot llegir, juntament amb la \FC, aproximadament un cop per segon.

\subsection{Intervals entre pujades}
Cada vegada que el comptador puja, tanca un interval $j$, entre els instants $t_{j-1}$ i
$t_j$, amb una energia
\begin{equation}
  \Delta E_j=E_{\tot}(t_j)-E_{\tot}(t_{j-1})=\int_{t_{j-1}}^{t_j}E(t)\,dt .
  \label{eq:deltaE}
\end{equation}
És la integral de l'energia en aquell interval, i no l'hem de calcular: la dona el comptador
de Garmin. Normalment val 1~kcal.

\paragraph{Els grams de cada interval.} Són la integral del ritme~\eqref{eq:cho-final} en
aquell interval:
\[
  \Delta\cho_j=\int_{t_{j-1}}^{t_j}\cho(t)\,dt
  =\int_{t_{j-1}}^{t_j}E(t)\cdot\frac{g_{\CHO}\bigl(\RER(t)\bigr)}{k\bigl(\RER(t)\bigr)}\,dt .
\]
L'interval dura uns 5~s, i hi prenem el \RER\ constant, igual a la mitjana dels \RER\ de cada
segon (figura~\ref{fig:comptador}):
\[
  \RERi_j=\frac{\sum_i\RER_i\cdot\Delta t_i}{\sum_i\Delta t_i} .
\]
Llavors $g_{\CHO}/k$ és un nombre fix i surt fora de la integral. La integral que queda és la
de l'energia, que és $\Delta E_j$~\eqref{eq:deltaE}:
\begin{align}
  \Delta\cho_j &\approx\frac{g_{\CHO}(\RERi_j)}{k(\RERi_j)}\cdot\int_{t_{j-1}}^{t_j}E(t)\,dt
  =\Delta E_j\cdot\frac{g_{\CHO}(\RERi_j)}{k(\RERi_j)}, \label{eq:interval}\\
  \Delta\fat_j &\approx\Delta E_j\cdot\frac{g_{\FAT}(\RERi_j)}{k(\RERi_j)} . \notag
\end{align}
L'única aproximació és prendre el \RER\ constant durant uns segons. Com que els intervals són
tan curts, gairebé no canvia el resultat.

\begin{figure}[H]
\centering
\begin{tikzpicture}
\begin{axis}[
  width=0.82\textwidth, height=4.6cm,
  xmin=0, xmax=24, ymin=99.4, ymax=104.7,
  xlabel={temps (s)}, ylabel={$E_{\tot}(t)$ (kcal)},
  ytick={100,101,102,103,104},
  tick label style={font=\small}, label style={font=\small},
  grid=major, grid style={gray!20}]
\addplot[const plot, very thick, black] coordinates {
 (0,100) (3,101) (8,102) (13,103) (18,104) (24,104)};
\draw[dashed, red!70!black, thick] (axis cs:1,99.4) -- (axis cs:1,104.7);
\node[font=\scriptsize, red!70!black, anchor=north west, align=left] at (axis cs:1.2,104.6) {primera\\lectura};
\draw[{Stealth}-{Stealth}, blue!70!black] (axis cs:1,100.5) -- (axis cs:3,100.5);
\node[font=\scriptsize, blue!70!black, anchor=south] at (axis cs:2,100.55) {parcial};
\draw[{Stealth}-{Stealth}, blue!70!black] (axis cs:3,101.5) -- (axis cs:8,101.5);
\node[font=\scriptsize, blue!70!black, anchor=south] at (axis cs:5.5,101.55) {interval $j$};
\draw[{Stealth}-{Stealth}, blue!70!black] (axis cs:8,102.5) -- (axis cs:13,102.5);
\node[font=\scriptsize, blue!70!black, anchor=south] at (axis cs:10.5,102.55) {interval $j+1$};
\end{axis}
\end{tikzpicture}
\caption{Comptador d'energia acumulada de Garmin. Cada pujada d'1~kcal tanca un interval,
l'energia del qual es reparteix amb~\eqref{eq:interval} i el \RER\ mitjà de l'interval. La
primera lectura cau en un punt qualsevol entre dues pujades, de manera que el primer
interval és parcial.}
\label{fig:comptador}
\end{figure}

\subsection{Totals}
Els totals~\eqref{eq:totals} són la integral de tota l'activitat, que és la suma de les
integrals de cada interval:
\begin{equation}
  \cho_{\tot}(T)=\int_0^T\cho(t)\,dt=\sum_j\int_{t_{j-1}}^{t_j}\cho(t)\,dt=\sum_j\Delta\cho_j ,
\end{equation}
i igual per al greix. Garmin fa la integral de l'energia per nosaltres; només cal sumar els
trossos.
Pel «principi de conservació de l'energia Garmin»~\eqref{eq:conservacio}, aplicat a cada
interval,
\begin{equation}
  e_{\CHO}\cdot\cho_{\tot}(T)+e_{\FAT}\cdot\fat_{\tot}(T)=E_{\tot}(T)-E_{\tot}(t_0):
\end{equation}
l'energia repartida és exactament la que ha comptat Garmin des de la primera lectura $t_0$,
sense cap error de suma. Si durant una estona no es llegeix el comptador, no es perd
energia: la pujada següent ja la inclou. L'energia que el comptador ja portava a la primera
lectura no es reparteix. Amb el cronòmetre aturat no s'acumula res.

% =====================================================================
\section{Exemple}\label{sec:exemple}
% =====================================================================

Un esportista amb $LT_2=166$~ppm i $\FCmax=180$~ppm corre a 150~ppm, i a partir del
comptador de Garmin el ritme de despesa és $E=12$~kcal/min.
\begin{enumerate}
  \item \textbf{\RER}~\eqref{eq:rer}: $LT_1=0{,}89\cdot166=147{,}7$~ppm, de manera que
        \[
          \RER=0{,}91+\frac{150-147{,}7}{166-147{,}7}\cdot(0{,}98-0{,}91)=0{,}919 .
        \]
  \item \textbf{Grams per litre}~\eqref{eq:gc}--\eqref{eq:gf}:
        $g_{\CHO}=4{,}210\cdot0{,}919-2{,}962=0{,}906$~g/L i
        $g_{\FAT}=1{,}695-1{,}701\cdot0{,}919=0{,}132$~g/L.
  \item \textbf{Energia per gram}~\eqref{eq:e}: $e_{\CHO}=4{,}1$ i $e_{\FAT}=9{,}3$~kcal/g.
  \item \textbf{Energia per litre}~\eqref{eq:k}:
        $k=4{,}1\cdot0{,}906+9{,}3\cdot0{,}132=4{,}944$~kcal/L.
  \item \textbf{Consum d'oxigen}~\eqref{eq:vo2}: $\VO=12/4{,}944=2{,}43$~L/min.
  \item \textbf{Ritme d'oxidació}: $\cho=2{,}43\cdot0{,}906=2{,}20$~g/min
        i $\fat=2{,}43\cdot0{,}132=0{,}32$~g/min, és a dir, 132~g/h de carbohidrats i
        19~g/h de greix.
\end{enumerate}
Condició de partida~\eqref{eq:conservacio}: $4{,}1\cdot2{,}198+9{,}3\cdot0{,}321=12{,}00$~kcal/min.

% =====================================================================
\section{Fonts d'error i limitacions}\label{sec:limitacions}
% =====================================================================

La taula~\ref{tab:errors} recull les fonts d'error, classificades segons si afecten la
quantitat d'energia o el seu repartiment.

\begin{table}[H]
\centering
\caption{Fonts d'error del mètode.}
\label{tab:errors}
\small
\begin{tabularx}{\textwidth}{@{}>{\raggedright\arraybackslash}p{3.6cm} L >{\raggedright\arraybackslash}p{2.4cm}@{}}
\toprule
Font & Descripció & Afecta \\
\midrule
Energia de Garmin & Càlcul del fabricant, amb errors coneguts en dispositius
portables \citep{shcherbina2017,fuller2020}; depèn del perfil de l'usuari (pes, edat,
\FC\ màxima, $\VOmax$). & Quantitat (i grams, en la mateixa proporció) \\
\RER\ estimat & Valors típics de població; cada persona pot variar uns $\pm0{,}04$
\citep{goedecke2000}; depèn dels llindars, sobretot de $LT_2$. & Repartiment \\
Durada de l'esforç & El \RER\ real baixa a mesura que avança l'esforç
\citep{rothschild2022}; el model no ho té en compte. & Repartiment (CHO sobreestimat a
partir de la segona hora) \\
Dieta, dejú, nivell & El dejú o una dieta baixa en carbohidrats baixen el \RER; vegeu
\citep{chofatfc2026}. & Repartiment \\
Proteïna & Les equacions no tenen en compte l'oxidació de proteïnes
\citep{frayn1983,jeukendrup2005}: la seva energia, petita, s'atribueix a carbohidrats i
greixos. & Repartiment (petit) \\
Glucòlisi anaeròbica & Per sobre de $LT_2$ es consumeix glicogen que no apareix a
l'oxidació calculada a partir dels gasos. & Repartiment (CHO subestimat) \\
\bottomrule
\end{tabularx}
\end{table}

El repartiment no es pot ajustar amb cap dada de Garmin, perquè Garmin no diu res sobre
carbohidrats i greixos. Per validar-lo cal una prova d'esforç amb analitzador de gasos.
L'energia, en canvi, és tan bona com l'estimació de Garmin: si aquesta millora, el resultat
millora sense canviar res.

% =====================================================================
\section{Estimació del consum mitjà de CHO en temps real}\label{sec:implementacio}
% =====================================================================

\paragraph{Què volem.} Una estimació en temps real de quants grams de carbohidrats per hora
s'estan oxidant, que doni més pes als últims minuts.

\paragraph{El problema.} Els grams arriben a salts. A cada pujada d'1~kcal (cada 5~s,
més o menys) entren de cop els grams de l'interval~\eqref{eq:interval}, i entre pujades no
n'entra cap. El ritme de cada segon és zero o un pic, i no serveix com a estimació.

\paragraph{Primera idea: la mitjana de tota l'activitat.} Dividir el total de grams pel
temps. És estable, però al cap d'una hora gairebé no es mou quan canvia la intensitat. Té
sentit com a resum al final de l'activitat.

\paragraph{La solució: una mitjana que pesa més els últims minuts.} Es fa servir una
mitjana mòbil exponencial \citep{hunter1986}. Amb paraules: a cada lectura (un cop per
segon), el valor anterior perd una petita part del seu pes, $\Delta t/\tau$, i hi entren
els grams nous dividits per $\tau$. Amb $\tau=15$~min:
\begin{equation}
  \bar r_i=\Bigl(1-\frac{\Delta t_i}{\tau}\Bigr)\cdot\bar r_{i-1}+\frac{\Delta\cho_i}{\tau},
  \label{eq:ema}
\end{equation}
on $\bar r_i$ és l'estimació, en g/min (multiplicada per 60, en g/h); $\Delta t_i$ és el
temps des de la lectura anterior, i $\Delta\cho_i$ els grams assignats en aquesta lectura
(zero si no hi ha hagut pujada).

\paragraph{Què vol dir $\tau=15$~min.} L'estimació és, a grans trets, el consum dels
\textbf{últims 15 minuts}, i el que és més recent hi pesa més. És com la ``memòria'' de la
mitjana: el que ha passat fa 15 minuts encara compta una mica, el de fa 30 minuts molt poc, i
el de fa 45 minuts pràcticament s'ha oblidat. Per exemple, si es passa d'un ritme suau a un
de fort, l'estimació no salta de cop: al cap de 15 minuts ja ha fet uns 2/3 del canvi, i al
cap de 45 minuts gairebé tot.

\paragraph{Per què 15 minuts.} Amb menys temps l'estimació reacciona més de pressa, però
balla més; amb més temps és més estable, però va endarrerida. 15 minuts és un compromís, i
és més o menys l'escala de temps en què es decideix quan menjar.

\paragraph{Inici.} El primer interval és parcial (figura~\ref{fig:comptador}), però porta
una quilocaloria sencera, de manera que el seu ritme pot sortir diverses vegades més alt
que el real. Si la mitjana comencés amb aquest valor, l'error tardaria molt a desaparèixer.
Per això, durant el primer minut es fa servir la mitjana real (grams dividits pel temps),
on aquest error ja queda repartit entre una dotzena d'intervals, i a partir d'aquí la mitjana
mòbil comença amb aquest valor.

% =====================================================================
\section{Conclusions}
% =====================================================================

El punt de partida d'aquest treball és una condició: l'energia dels grams calculats ha de
ser la que dona Garmin, que fa servir més informació que la \FC\ sola. El mètode no torna a
calcular l'energia; només la reparteix entre carbohidrats i greix.

Les equacions de Jeukendrup i Wallis calculen el ritme d'oxidació de carbohidrats i greix
a partir del consum d'oxigen i del \RER. Sense analitzador de gasos, el \RER\ es pot
estimar amb la \FC, i el consum d'oxigen amb l'energia que dona Garmin: n'hi ha prou de
dividir-la per l'energia que dona cada litre d'oxigen, $k(\RER)$. Aquesta peça clau es
calcula amb les mateixes equacions i amb l'energia de cada gram de carbohidrat (4,1~kcal)
i de greix (9,3~kcal), uns valors que encaixen amb la literatura i amb la calorimetria
clàssica. Com que $k$ varia poc, el consum d'oxigen surt gairebé només de l'energia. I com
que es calcula amb els mateixos nombres, l'energia dels grams obtinguts és exactament la de
Garmin, tant en cada instant com en el total. Al final hi ha dues fonts d'error ben
separades: la quantitat d'energia, que ve de Garmin, i el repartiment entre carbohidrats i
greix, que depèn del \RER\ estimat i, sobretot, de si els llindars estan ben posats.

El mètode serviria amb qualsevol dispositiu que doni l'energia total de l'activitat, però
les comprovacions sobre com es comporta el comptador d'energia s'han fet amb Garmin.

% =====================================================================
\begin{thebibliography}{99}
\small

\bibitem[Atwater i Bryant(1900)]{atwater1900}
Atwater WO, Bryant AP.
\newblock The availability and fuel value of food materials.
\newblock \emph{Storrs Agricultural Experiment Station 12th Annual Report}. 1900:73--110.

\bibitem[Firstbeat Technologies(2007)]{firstbeat2007}
Firstbeat Technologies Ltd.
\newblock An energy expenditure estimation method based on heart rate measurement.
\newblock White paper. Jyväskylä: Firstbeat Technologies; 2007.

\bibitem[Frayn(1983)]{frayn1983}
Frayn KN.
\newblock Calculation of substrate oxidation rates in vivo from gaseous exchange.
\newblock \emph{J Appl Physiol}. 1983;55(2):628--634.

\bibitem[Fuller et~al.(2020)]{fuller2020}
Fuller D, Colwell E, Low J, et~al.
\newblock Reliability and validity of commercially available wearable devices for
  measuring steps, energy expenditure, and heart rate: systematic review.
\newblock \emph{JMIR Mhealth Uhealth}. 2020;8(9):e18694.

\bibitem[Garmin(2026)]{garminciq}
Garmin Ltd.
\newblock Connect IQ SDK 9.2.0, documentació de l'API: \texttt{Toybox.Activity.Info}.
  2026.

\bibitem[Goedecke et~al.(2000)]{goedecke2000}
Goedecke JH, St~Clair Gibson A, Grobler L, Collins M, Noakes TD, Lambert EV.
\newblock Determinants of the variability in respiratory exchange ratio at rest and during
  exercise in trained athletes.
\newblock \emph{Am J Physiol Endocrinol Metab}. 2000;279(6):E1325--E1334.

\bibitem[Hunter(1986)]{hunter1986}
Hunter JS.
\newblock The exponentially weighted moving average.
\newblock \emph{J Qual Technol}. 1986;18(4):203--210.


\bibitem[Jeukendrup(2014)]{jeukendrup2014}
Jeukendrup AE.
\newblock A step towards personalized sports nutrition: carbohydrate intake during
  exercise.
\newblock \emph{Sports Med}. 2014;44(Suppl 1):S25--S33.

\bibitem[Jeukendrup i Wallis(2005)]{jeukendrup2005}
Jeukendrup AE, Wallis GA.
\newblock Measurement of substrate oxidation during exercise by means of gas
  exchange measurements.
\newblock \emph{Int J Sports Med}. 2005;26(Suppl 1):S28--S37.

\bibitem[Livesey i Elia(1988)]{livesey1988}
Livesey G, Elia M.
\newblock Estimation of energy expenditure, net carbohydrate utilization, and net fat
  oxidation and synthesis by indirect calorimetry: evaluation of errors with special
  reference to the detailed composition of fuels.
\newblock \emph{Am J Clin Nutr}. 1988;47(4):608--628.

\bibitem[Péronnet i Massicotte(1991)]{peronnet1991}
Péronnet F, Massicotte D.
\newblock Table of nonprotein respiratory quotient: an update.
\newblock \emph{Can J Sport Sci}. 1991;16(1):23--29.

\bibitem[Projecte CHO/FAT Metabolic Pro(2026)]{chofatfc2026}
Projecte CHO/FAT Metabolic Pro.
\newblock Estimació de l'oxidació de carbohidrats i greixos a partir de la freqüència
  cardíaca. 2026.
\newblock \webprojecte

\bibitem[Rothschild et~al.(2022)]{rothschild2022}
Rothschild JA, Kilding AE, Stewart T, Plews DJ.
\newblock Factors influencing substrate oxidation during submaximal cycling: a modelling
  analysis.
\newblock \emph{Sports Med}. 2022;52(11):2775--2795.

\bibitem[Shcherbina et~al.(2017)]{shcherbina2017}
Shcherbina A, Mattsson CM, Waggott D, et~al.
\newblock Accuracy in wrist-worn, sensor-based measurements of heart rate and energy
  expenditure in a diverse cohort.
\newblock \emph{J Pers Med}. 2017;7(2):3.

\bibitem[Swain i Leutholtz(1997)]{swain1997}
Swain DP, Leutholtz BC.
\newblock Heart rate reserve is equivalent to \%VO$_2$ reserve, not to
  \%VO$_2$max.
\newblock \emph{Med Sci Sports Exerc}. 1997;29(3):410--414.

\bibitem[Weir(1949)]{weir1949}
Weir JB de V.
\newblock New methods for calculating metabolic rate with special reference to
  protein metabolism.
\newblock \emph{J Physiol}. 1949;109(1--2):1--9.

\end{thebibliography}

\end{document}
